% Estilo común de los esquemáticos (circuitikz)
\usepackage[siunitx, european resistors, american inductors, straightvoltages]{circuitikz}
\usepackage{siunitx}
\sisetup{output-decimal-marker={,}}
\ctikzset{bipoles/length=0.95cm, diodes/scale=0.85, resistors/scale=0.85, capacitors/scale=0.85}
\tikzset{font=\footnotesize, lbl/.style={font=\footnotesize, inner sep=1pt}, net/.style={font=\footnotesize\bfseries, inner sep=1pt}}
% Regulador de tres terminales: \reg{(x,y) del pin IN}{nombre}{pin IN}{pin OUT}{pin ADJ}; caja de 2,6 x 1,4
\newcommand{\reg}[5]{%
  \draw[thick] ($#1+(0,-0.7)$) rectangle ++(2.6,1.4);
  \node[font=\footnotesize\bfseries] at ($#1+(1.3,0.32)$) {#2};
  \node[font=\tiny, anchor=west] at ($#1+(0.05,-0.15)$) {IN};
  \node[font=\tiny, anchor=east] at ($#1+(2.55,-0.15)$) {OUT};
  \node[font=\tiny, anchor=south] at ($#1+(1.3,-0.68)$) {ADJ};
  \node[font=\tiny, anchor=south east] at ($#1+(-0.02,0.02)$) {#3};
  \node[font=\tiny, anchor=south west] at ($#1+(2.62,0.02)$) {#4};
  \node[font=\tiny, anchor=north west] at ($#1+(1.33,-0.72)$) {#5};
}
% etiqueta apilada: nombre arriba, valor abajo
\newcommand{\lab}[2]{\shortstack[l]{#1\\#2}}
